\documentclass[aspectratio=43,8pt]{beamer}
\usepackage[utf8]{inputenc}
\usepackage{multicol}
\usepackage{hyperref}
\usepackage{babel}[french]
\hypersetup{
    colorlinks=true,
    urlcolor=blue,
    pdfborder={0 0 0},  % pas de cadre moche
}

\title{On Demand Buffer (ODB) - Application à Nginx}
\author[Brelot]{Julien Brelot}
\date{ Work in Progress , \today}

\usetheme{material}
\useLightTheme
\usePrimaryGreen
\useAccentCyan

\setbeamertemplate{footline}[text line]{%
  \hfill%
  \insertframenumber\,/\,\inserttotalframenumber
}

\begin{document}

\begin{frame}
\titlepage
\end{frame}

\begin{frame}{Nginx : Principe de fonctionnement}
\begin{card}
    Pour traiter les requêtes, nginx utilise différents processus :
    \begin{itemize}
        \item   Master : Charge la configuration et se charge de créer la socket d'écoute du serveur. Il crée et manage les autres processus.  
        \item   Cache Loader : charge la mémoire cache afin de répondre plus rapidement aux clients accédant à divers ressources.
        \item   Cache Manager : régule la mémoire cache quand la limite a été atteinte.
        \item Worker : Accepte les connexions, gère les requêtes, les traite et y répond efficacement grâce à sa boucle événementiel.
    \end{itemize}
\end{card}
\end{frame}

\begin{frame}{Nginx : Boucle de gestion d'évènements}
\begin{card}
    Afin de gérer efficacement un grand nombre de connexion,le worker nginx se repose sur un mécanisme de gestion d'évènement :
    \begin{itemize}
        \item   Chaque descripteur de fichier est ajouté dans une file d'observations.
        \item   Lorsqu'un fichier (ou ensemble de fichiers) est disponible (en lecture, ecriture,...), il est ajoutré dans une queue d'évènement à traiter.
        \item   Chaque évènement de la queue est ensuite traité.
        \item   Une fois traité, le descripteur associé est retiré de la queue.
        \item   On réitère jusquà ce que la queue soit vide.
        \item   Lorsque la queue est vide, le worker attend un nouvel évènement.
    \end{itemize}
	Le worker enregistre donc une liste de descripteurs à surveiller et attend qu'un évènement se produise sur au moins l'un d'eux.
	Tout cela s'effectue séquentiellement.\\
	Ainsi, nginx repose donc sur l'utilisation d'un mécanisme d'obersvation et d'attentes d'évènements comme poll (POSIX), epoll (Linux) ou kqueue (FreeBSD).\\
	Il est donc nécesssaire que les sockets soient non bloquantes afin de réaliser les opérations de lecture/ écritures quand elles sont disponibles uniquement.
\end{card}
\end{frame}

\begin{frame}{Nginx : Boucle de gestion d'évènements}
\begin{figure}
        \centering
        \cardImg{img/NGINX-Event-Loop.png}{0.8\textwidth}
        \caption{\small Nginx Event Loop (source: \href{https://www.f5.com/company/blog/nginx/thread-pools-boost-performance-9x}{https://www.f5.com})}
\end{figure}

\end{frame}

\begin{frame}{Nginx : Gestion d'un évenèment}
\centering
\begin{figure}
        \centering
        \cardImg{img/NGINX-Events-Queue-Processing-Cycle.png}{0.8\textwidth}
        \caption{\small Event Cycle (source: \href{https://www.f5.com/company/blog/nginx/thread-pools-boost-performance-9x}{https://www.f5.com})}
\end{figure}
\end{frame}

\begin{frame}{Nginx : Opération lente (I/O intensive)}
\begin{card}
	Les opérations I/O peubvent être très lentes et ralentir la boucle de gestion d'évènements.\\
	Pour cela, nginx a mis en place des thread pools afin de gérer les opérations I/O lourdes en parallèle et éviter de bloquer un worker durant ces opérations.
\end{card}
\begin{figure}
        \centering
        \cardImg{img/nginx-mode.png}{0.8\textwidth}
        \caption{\small epoll mode (source: \href{https://medium.com/@avocadi/what-is-epoll-9bbc74272f7c}{https://medium.com})}
\end{figure}

\end{frame}

\begin{frame}{Nginx : Opération lente (I/O intensive)}
\begin{figure}
        \centering
        \cardImg{img/nginx-thread-pool.png}{0.8\textwidth}
        \caption{\small Nginx Thread Pool (source: \href{https://www.f5.com/company/blog/nginx/thread-pools-boost-performance-9x}{https://www.f5.com})}
\end{figure}
\end{frame}

\begin{frame}{Le cas Linux : epoll}
\begin{card}
	Sur Linux, nginx utilise epoll pour surveiller les sockets.\\
	En,effet, dans le cas d'un grand nombre de sockets, epoll est plus performant que poll.\\
	Il existe 2 modes d'utilisation de epoll :
	\begin{itemize}
		\item   Epoll Edge Triggered
		\item   Epoll Level Triggered
	\end{itemize}
	Le mode Edge Triggered est celui utilisé par nginx.\\
	Ce mode permet de déclencher un évènement lorsque l'état du fichier change.\\
	A contrario, le mode Level Triggered permet de déclencher un évènement lorsque le fichier est disponible (lecture, ecriture,...).\\
\end{card}
\end{frame}

\begin{frame}{Problème : epoll Edge Triggered}
\begin{card}
	Le mode Edge Triggered est celui utilisé par nginx.\\
	De par son principe, il peut générer des soucis de famines ou dans notre cas, amener des données à ne pas être lues. \\
	Pour illustrer ce problème, prenons cet exemple :
	\begin{itemize}
		\item   Un back-end (BE) envoie 2Ko de données au serveur nginx.
		\item   Supposons, le buffer de nginx est de 1Ko.
		\item   Le BE envoie 2Ko de données au serveur nginx.
		\item 	A la réception des données, un event est déclenché (changement d'état de  la socket).
		\item   La socket devient lisible, nginx lit 1Ko de données.
		\item 	Il reste 1 Ko dans le buffer du noyau.
		\item   Or si nginx ne vérifie pas qu'il a bien recu toutes les données, il devra se contenter de 1Ko de données.
		\item   Pourquoi ? La socket reste dans le même etat (lisible), pas de nouvel event reçu.
	\end{itemize}
\end{card}
\end{frame}

\begin{frame}{Problème : epoll Edge Triggered(2)}
\begin{card}
	En l'occurence, nginx s'assure d'avoir bien lues toutes les données jusqu'à EAGAIN.\\
	Or ODB, fait croire, dans le cas de données virtuelles, qu'il lu toutes les données demandées par recv.\\
	De ce fait, il peut rester d'autres données dans le buffer de réception du noyau.\\
	Ces données peuvent être d'autres descripteurs virtuels .\\
	Ainsi, nginx peut ne pas lire toutes les données virtuelles.\\

	\textbf{Solution 1 :} Détecter à la fin du recv, si le buffer du noyau est vide. Si non, générer utiliser  epoll\_ctl pour signaler un event à nginx.\\
	\textbf{Solution 2 :} Intercepter epoll\_ctl (enrgistre l'intérêt pour un descripteur), pour utiliser le mode Level Triggered.\\

    La solution 1 a été explorée et semble fonctionnée. Il a fallu hooké epoll\_create et epoll\_ctl pour récupérer le descripteur vers le epoll et les données que nginx passe en paramètre (des callbacks) aux events (struct epoll\_event).
\end{card}
\end{frame}

\begin{frame}{Ressources utilisées}
\begin{card}
    Ressources sur nginx :
    	\begin{itemize}
    		\item \url{https://nginx.org/en/docs/dev/development_guide.html#event_loop}
    		\item \url{https://medium.com/%40rahuls163284/the-magic-of-nginx-understanding-nginxs-event-driven-system-27866d6e72bb}
    		\item \url{https://codedamn.com/news/backend/nginx-connection-processing}
    		\item \url{https://medium.com/@soulaimaneyh/nginx-internal-architecture-b94b013bc365}
    		\item \url{https://youtu.be/I6dpN0geIb4}
    		\item \url{https://www.f5.com/company/blog/nginx/thread-pools-boost-performance-9x}
    	\end{itemize}
    Ressources sur epoll :
        \begin{itemize} 
        	\item \url{https://medium.com/@avocadi/what-is-epoll-9bbc74272f7c}
        	\item \url{https://man.archlinux.org/man/epoll.7.fr}
        \end{itemize}
\end{card}
\end{frame}

\end{document}